\paragraph{Formal relationship to neighboring methods.} The distinction is between what each observation constrains and what further assumptions permit inference.
\begin{table}[htbp]\centering\small
\caption{Measurement constraints and the additional inference made here.}\label{tab:formal-comparison}
\begin{tabularx}{\linewidth}{@{}p{.24\linewidth}YY@{}}\toprule
Approach & What it constrains & Additional requirement\\\midrule
Phase-slope index \citep{nolte2008} & Frequency-dependent cross-spectral phase and a directional summary & A physiological delay interpretation requires constrained transfer and processing phases.\\\addlinespace
Delay/interference analysis \citep{dowdall2023} & Effects of delays on coherence and Granger estimates in interacting systems & The present ridge concerns an explicit conditional carrier/gate likelihood, not the same interference mechanism.\\\addlinespace
Communication subspaces \citep{semedo2019} & Predictive population directions across areas & A causal pathway and a threshold/readout rule require independent calibration and interventions.\\\addlinespace
This model & Offset equivalence, projected item contrasts, return products and operational first passage & Known observation channels, constrained phase branches, calibrated maps, and module-specific tests.\\\bottomrule
\end{tabularx}\end{table}
Granger causality assesses whether past signals improve prediction within a specified multivariate time-series model. Broadband lag structure can add information beyond the stationary one-band carrier observation; the ridge does not prove that all Granger analyses are uninformative. Conversely, a selected predictive lag is not automatically an anatomical conduction time, particularly under feedback, unknown filtering, or omitted inputs. The empirical comparison should fit those rival time-series explanations to the same held-out perturbation responses, rather than infer a delay from a single summary.
